\documentclass[runningheads]{llncs}
\usepackage[T1]{fontenc}
\usepackage{amsmath,amssymb,booktabs,array,multirow}
\usepackage{graphicx}
\usepackage{float}
\usepackage{xurl}
\usepackage{microtype}
\usepackage[hidelinks]{hyperref}
\newcommand{\ECL}{\mathrm{ECL}}
\newcommand{\PD}{\mathrm{PD}}
\newcommand{\EAD}{\mathrm{EAD}}
\newcommand{\LGD}{\mathrm{LGD}}
\newcommand{\paperfigure}[2]{\includegraphics[width=#2]{../figures/#1}}
\title{The Quantity--Risk Trade-off in Feedback Acquisition under Selective Labels}
\titlerunning{Risk-Budgeted Feedback Acquisition under Selective Labels}
\author{Anonymous Authors}
\authorrunning{Anonymous Authors}
\institute{Anonymous Institution}
\begin{document}
\maketitle
\begin{abstract}
Lending decisions determine which repayment outcomes enter future training data. When feedback from a rejected region requires approving applicants, an acquisition score changes both the number of observed outcomes and the exposure used to obtain them. We define a maximum-feasible-count reference, $K^*$, and a same-count minimum-risk reference $M(k)$, then measure the quantity opportunity cost and predicted-risk premium of score-based selection. Random-order replay covers four credit datasets and five budget ratios; a registered five-policy block compares FAVE with bootstrap-committee information gain (IG). At $\rho=0.1$, LRBE-IG acquires 31.3, 49.5, and 40.1 more labels than ScoreFirst-IG on HMEQ, LendingClub, and UCI. Relative to MinRisk, it uses 28,765.5 and 27,025.7 additional predicted-risk units on HMEQ and LendingClub, while Recall and Brier changes vary by dataset and direction. Same-state audits and a controlled retraining experiment connect these policy outcomes to immediate selection choices. The evidence supports evaluating acquisition scores by feedback quantity, predicted-risk use, composition, and downstream performance together.
\keywords{Selective labels \and Feedback acquisition \and Risk budgets \and Credit prediction}
\end{abstract}
\section{Introduction}
\label{sec:introduction}

Credit decisions determine which data become available for subsequent learning. Repayment outcomes are observed for approved applicants, while the outcomes of rejected applications remain unknown. This selective-label problem can leave evaluation and retraining dependent on the policy that collected the data \cite{lakkaraju2017selective,kleinberg2018human}. Acquiring feedback from a previously rejected region therefore requires a lending decision and acceptance of the associated credit exposure.

Existing research addresses this dependence through reject inference \cite{banasik2007reject}, decision policies for selective observations \cite{wei2021selective}, optimistic exploration \cite{pacchiano2021plot}, and adaptive bounded debiasing \cite{yang2022bounded}. Recent credit experiments also use controlled approvals to investigate feedback loops \cite{scarone2026illusion}. These approaches motivate a practical question about resource allocation: under a predicted-loss budget, does sacrificing feasible feedback quantity for a higher information score deliver sufficient downstream benefit? Non-uniform acquisition costs are already studied in active learning \cite{donmez2008proactive,tomanek2010cost}; our focus is the comparison of quantity-preserving and composition-oriented policies when the cost is an applicant-level credit-loss proxy.

We examine this question using lexicographic risk-budgeted feedback acquisition (LRBE). At each round, the first stage computes the maximum feasible count $K_t^*$ under predicted costs $\widehat c_{it}=\widehat p_{it}\mathrm{EAD}_i\mathrm{LGD}_i$. The second stage optimizes a sample-composition score subject to that count equality. A secondary score can therefore change which applicants are selected without reducing the current maximum count. Removing the equality produces a composition-first policy and exposes the number of feasible observations exchanged for a higher score. Minimum-risk, joint-objective, and randomized policies provide complementary references.

The evaluation combines three views of the trade-off. Dynamic replay on four credit datasets measures complete policy outcomes over five budget ratios. A five-policy experiment at $\rho=0.1$ compares FAVE with bootstrap-committee information gain (IG), whose score is independent of FAVE's Fisher geometry and influence approximation. Same-state audits then give every selector the same visible set, candidates, costs, and numerical budget, while a controlled one-step experiment varies exposure heterogeneity and the evaluation weight of an underrepresented region. Dynamic policies share their initial data and incoming batches but subsequently retrain on their own feedback, so later candidate sets and budgets may differ.

The experiments reveal an asymmetric trade-off. Composition-first FAVE and IG often give up many feasible observations, while their predictive estimates vary across datasets. LRBE-IG acquires 31.3, 49.5, and 40.1 more labels than QualityFirst-IG on HMEQ, LendingClub, and UCI, respectively; its Recall intervals span both directions and its UCI Brier error increases. Against MinRisk, FAVE and IG achieve similar cumulative counts while consuming more predicted risk, again with mixed predictive estimates. The controlled experiment reproduces the immediate count and risk differences, while its prespecified Brier interactions remain uncertain.

This study provides an empirical framework for evaluating risk-budgeted acquisition. The LRBE comparison separates the maximum feasible count at the current state from the score used to choose applicants, exposing the number of observations exchanged for composition. Applying the same construction to FAVE and committee IG, with randomized controls, measures the associated changes in feedback quantity, predicted-risk use, Recall, and Brier score. Same-state analyses identify the immediate selection trade-off, and the controlled experiment tests exposure heterogeneity and target-region emphasis as possible moderators of predictive return. Together, these analyses turn a broad appeal to ``informative feedback'' into a measurable comparison with minimum-risk maximum-count selection.

\section{Related Work}
\label{sec:related}

\paragraph{Learning from selectively observed outcomes.}
Work on selective labels studies the evaluation and learning problems created by decision-dependent observation \cite{lakkaraju2017selective,kleinberg2018human}. Expert consistency and stochastic fair decisions provide ways to learn or collect data beyond existing acceptance patterns \cite{dearteaga2018expert,kilbertus2020fair}. Credit reject inference examines augmentation and sample-selection correction \cite{banasik2007reject,marshall2010selection}, while finite-domain policy optimization and one-sided-feedback learning formalize sequential decisions with missing outcomes \cite{wei2021selective,jiang2021onesided}. Our replay retains observed historical outcomes and reveals them according to simulated approvals. Its estimand is acquisition-policy behavior under replay, distinct from outcome imputation and field experimentation on actual rejected borrowers.

\paragraph{Exploration objectives and costs.}
PLOT uses optimistic pseudo-labels to counter self-confirming rejection and gives regret guarantees under stated modeling assumptions \cite{pacchiano2021plot}. Adaptive bounded exploration trades debiasing against the risk of accepting otherwise rejected individuals \cite{yang2022bounded}. Partial-feedback exploration also supports subgroup coverage and fairness objectives \cite{keswani2024partial}, while cost-efficient fairness auditing accounts for assessment, processing, and potential-default costs \cite{das2025cost}. Controlled approvals in credit scoring examine the consequences of retraining under selection bias \cite{scarone2026illusion}. These works establish the surrounding problem of costly exploration. We compare how a current predicted-loss allowance is allocated between feedback quantity and composition, using randomized policies as strategy controls for the quantity constraint and selection score.

\paragraph{Resource constraints and predictive evaluation.}
Cost-sensitive active learning already accounts for heterogeneous query costs \cite{donmez2008proactive,tomanek2010cost}. Contextual bandits with knapsack constraints and conservative exploration place learning within resource or baseline-performance restrictions \cite{agrawal2016knapsacks,wu2016conservative}; safety-constrained approaches further consider uncertainty in feasibility \cite{amani2019safe,pacchiano2021linearconstraints}. LRBE uses a point predicted-loss constraint in place of these sequential guarantees. We evaluate secondary scores through the resulting data and predictions, comparing Recall and Brier while reporting predicted-budget compliance separately from realized defaulted exposure.

\section{Risk-Budgeted Feedback Acquisition}
\label{sec:problem}\label{sec:method}

\subsection{Sequential observation and policy-specific budgets}
At round $t$, policy $m$ receives an incoming batch $\mathcal U_t$ and fits a default predictor on its visible labeled set $\mathcal L_t^{(m)}$. An outcome $y_i=1$ denotes default. Ordinary approvals have predicted default probability below the fixed threshold $\tau=0.5$, which serves as a common operating point for every policy in the replay. The remaining applicants form the exploration candidates
\begin{equation}
\mathcal C_t^{(m)}=\{i\in\mathcal U_t:\widehat p_{it}^{(m)}\ge\tau\}.
\end{equation}
The policy selects exploratory approvals $\mathcal S_t^{(m)}\subseteq\mathcal C_t^{(m)}$. Only outcomes of ordinary and exploratory approvals are added to the next visible set. Rejected applicants are not carried over to a backlog. Outcomes are revealed within the simulation round; actual repayment delays are not modeled.

Write $a_i$ for exposure and set $\LGD_i=1$. The planning cost, exploration budget, and selected risk are
\begin{equation}
\widehat c_{it}^{(m)}=\widehat p_{it}^{(m)}a_i,\qquad
B_t^{(m)}=\rho\sum_{i\in\mathcal C_t^{(m)}}\widehat c_{it}^{(m)},\qquad
\widehat R_t^{(m)}=\sum_{i\in\mathcal S_t^{(m)}}\widehat c_{it}^{(m)}.
\label{eq:budget}
\end{equation}
The constraint $\widehat R_t^{(m)}\le B_t^{(m)}$ covers exploratory approvals only, not the entire lending portfolio. Defaulted exposure is evaluated separately as $R_t^{\mathrm{real},(m)}=\sum_{i\in\mathcal S_t^{(m)}}y_i a_i$. With unit LGD this is an experimental loss proxy, not recovery-adjusted net loss or full IFRS~9 accounting.

Policies share the initial split and exogenous batches, but their feedback choices affect subsequent training sets, predictions, candidates and budgets. Pairing on dataset, seed and $\rho$ therefore compares complete policies under common incoming data. It does not fix their subsequent numerical budgets or isolate a causal effect of composition. Superscripts $m$ are suppressed below for within-policy optimization.

\subsection{Lexicographic quantity and composition}
LRBE first determines the maximum feasible exploratory count:
\begin{equation}
K_t^*=\max_{z\in\{0,1\}^{|\mathcal C_t|}}\sum_i z_i
\quad\text{subject to}\quad \sum_i z_i\widehat c_{it}\le B_t.
\label{eq:kstar}
\end{equation}
For nonnegative costs and unit count value, sorting costs and taking the longest feasible prefix solves this stage in $O(n\log n)$ time. This elementary cardinality construction is used as a reference, not claimed as a new optimization algorithm. A fixed, outcome-blind secondary score $q_i$ then determines the composition:
\begin{equation}
\max_z\sum_i z_iq_i
\quad\text{subject to}\quad \sum_i z_i=K_t^*,\quad
\sum_i z_i\widehat c_{it}\le B_t,\quad z_i\in\{0,1\}.
\label{eq:stage2}
\end{equation}
The equality prevents the secondary objective from sacrificing feasible count \emph{at the current policy state}. It does not force distinct policies to acquire identical cumulative counts after their histories diverge. Both stages use only currently available features and predictions, never unobserved candidate outcomes or test labels.

\subsection{Acquisition policies}
\textbf{CountOnly-MinRisk} minimizes predicted risk among maximum-count sets by setting $q_i=-\widehat c_i$. It is abbreviated MinRisk in tables. \textbf{LRBE-Uncertainty} uses $q_i=\widehat p_i(1-\widehat p_i)$. \textbf{LRBE-FAVE} uses the frozen FAVE-v2.1 score, a feature-space local information proxy (Appendix~\ref{app:fave}). Its Fisher-style geometry is a surrogate constructed from predictions and covariates, not the Hessian of the fitted boosted-tree model or a guarantee of retraining improvement.

\textbf{LRBE-IG} uses empirical committee mutual information. At each round, 20 models are fitted to same-size bootstrap samples of the current visible set. If $p_i^{(r)}$ is member $r$'s default probability and $H(p)=-p\log p-(1-p)\log(1-p)$, then
\begin{equation}
q_i^{\mathrm{IG}}=H\!\left(\frac{1}{20}\sum_{r=1}^{20}p_i^{(r)}\right)-\frac{1}{20}\sum_{r=1}^{20}H\!\left(p_i^{(r)}\right).
\label{eq:ig}
\end{equation}
Committee members use the canonical XGBoost hyperparameters and only visible labels. The committee determines acquisition scores only; the single canonical model continues to define candidates, predicted costs, and final evaluation. Scores add across selected applicants, so this comparator is a single-item IG sum rather than joint BatchBALD. \textbf{QualityFirst-IG} maximizes the same nonnegative score under the risk constraint without the count equality. IG is an independent acquisition baseline, not a claimed contribution.

\textbf{Quality-first policies} maximize the corresponding nonnegative uncertainty or FAVE utility under the risk constraint without the count equality. \textbf{Joint-Best-Val-Uncertainty} maximizes $\sum_i z_i(\bar q_i-\lambda\bar c_i)$ over positive-score candidates, where bars denote within-pool min--max normalization. Its parameter is chosen using only the initial visible set, split into 60\% training, 20\% pseudo-candidates and 20\% validation. Validation PR-AUC determines the choice, with Recall and then smaller $\lambda$ as tie-breakers. These are diagnostic comparators for the acquisition objectives, not implementations of all prior exploration algorithms.

Two additional controls vary randomization. \textbf{RiskBounded-Random} visits candidates in a reproducible pseudo-random order and accepts each if its cost fits the remaining budget. \textbf{RiskBounded-Epsilon} instead preserves $K_t^*$ and uses $q_i=-0.9\bar c_i+0.1u_i$, where $u_i$ is a deterministic hash-based random score. We retain its archive name but call it the \emph{random-score} control in prose: 0.1 is an objective weight, not a 10\% random-approval probability or conventional epsilon-greedy rule. \textbf{NoExplore} reveals ordinary approvals only. \textbf{Oracle} also reveals every candidate outcome and ignores the exploration budget; it is a full-feedback diagnostic, not a guaranteed upper bound on test performance.

\section{Experimental Design}
\label{sec:experiments}
\subsection{Data and replay protocol}
\begin{table}[t]
\caption{Exact input cohorts and exposure semantics. Raw risk units are not comparable across datasets.}
\label{tab:datasets}
\centering
\begin{tabular}{lrl}
\toprule
Dataset & Records & Exposure and evidence role\\
\midrule
HMEQ & 5,960 & LOAN amount; development\\
LendingClub & 16,566 & loan-amount proxy; external validation\\
UCI & 30,000 & unit exposure; development\\
PPDai & 55,596 & unit exposure; exploratory external\\
\bottomrule
\end{tabular}
\end{table}
For each seed, a random permutation assigns 25\% of records to a fixed outer test set. The first 20\% of the remaining stream form the initial visible set, with a minimum of 100 records; all remaining records are split into ten incoming batches. The initial visible set is therefore a random cohort sample rather than the output of a historical approval rule. Policy-dependent selective observation begins only when the policies process the ten subsequent batches: ordinary approvals are observed automatically, whereas exploratory approvals are chosen from the threshold-defined candidate set. Thus the experiment is \emph{random-order offline replay}, not historical-time evaluation.

Fresh confirmation seeds 1001--1010 are disjoint from development seeds 0--19. The fresh design includes an objective comparison at $\rho=0.1$, a four-policy sweep over $\rho\in\{0.025,0.05,0.1,0.2,0.4\}$, and NoExplore/Oracle references. The two randomization controls use the same fresh seeds at $\rho=0.1$. A separate primary IG block at $\rho=0.1$ reruns MinRisk, LRBE-FAVE, QualityFirst-FAVE, LRBE-IG, and QualityFirst-IG together, yielding 200 dynamic episodes. Historical control outputs are not spliced into that block. Development analyses are retained separately because their protocols differ from the fresh uncertainty estimates. All reported analyses use the registered seeds, policies, and archived raw outputs.

A controlled one-step experiment uses synthetic tasks with prespecified seeds 2001--2020. Each task contains a selectively observed initial set, one incoming batch, and separate test strata for an initially underrepresented region and its complement. It crosses unit exposure with normalized lognormal exposure and evaluates region-1 weights of 0.1 and 0.5. MinRisk, LRBE-IG, QualityFirst-IG, and a maximum-count random-score selector share the fitted pre-acquisition state, candidates, costs, numerical budget, and branch-training seed. The experiment measures immediate count sacrifice and predicted-risk use, followed by one retraining step; it tests specified condition interactions rather than long-run policy behavior.

\subsection{Model, optimization and observation controls}
All fresh policies use XGBoost with 100 trees, maximum depth 3, learning rate 0.1, binary logistic objective, and histogram-based training. Numeric features receive median imputation and standardization; categorical features receive a missing-value category and one-hot encoding. Target, ID and the designated exposure column are excluded from predictors. Preprocessing is fitted on visible records only. Each round fits a pre-acquisition model for selection and a post-acquisition model for test evaluation. The outer test set never selects parameters or candidates; population default rates and selected outcomes are used only in offline diagnostics.

Maximum-count selection uses a sorted-cost prefix with a $10^{-10}$ numerical allowance. The common second-stage binary program and quality-first/joint programs use SciPy's mixed-integer solver; an unsuccessful solve raises an error. Risk auditing permits a $10^{-6}$ absolute allowance. Candidate costs are strictly positive under the validated exposures and rejection threshold, so zero budgets yield zero feasible count. Empty candidate sets also yield no exploratory selections. GPU training and supported FAVE matrix operations accelerate computation, while optimization remains CPU-side. The registered bootstrap committee runs on CPU with four concurrent workers; the canonical model remains on GPU. Appendix~\ref{app:repro} records implementation and validation details.

\subsection{Outcomes and statistical comparisons}
We report cumulative exploratory count, cumulative predicted risk, final-test default Recall, and Brier score. Recall is $TP/(TP+FN)$ at $\tau=0.5$; Brier is the mean squared error of default probabilities. The acquired default fraction describes the selected outcomes. A separate feedback diagnostic, $F$, is the default fraction in the entire visible set divided by the default fraction in the non-test stream. It therefore includes ordinary approvals and the initial set, not only exploratory selections. Neither $F$ nor mean selected predicted probability establishes exact applicant overlap or an independent composition effect.

At each dataset and budget ratio, final policy differences are paired over seeds. Tables show mean differences and pointwise 95\% percentile bootstrap intervals from 10,000 seed resamples. Counts that are close descriptively are not called experimentally matched; intervals crossing zero are not equivalence or noninferiority tests. The comparisons are exploratory across multiple metrics and conditions, so isolated interval exclusions are not familywise significance claims.

The IG block adds a same-state mechanism audit. At each round of the MinRisk trajectory, all five selectors receive the same visible set, candidates, cost vector, and numerical budget. We record count sacrifice $D(S)=K_t^*-|S|$, composition premium $P(S)=\widehat R(S)-M(|S|)$, where $M(k)$ is the minimum predicted risk of a size-$k$ set, maximum-count slack, selected IDs, and Jaccard overlap with MinRisk. This audit describes immediate selector behavior; because alternative selections do not update the anchor trajectory, it is not a downstream counterfactual.

The controlled experiment applies the same 10,000-draw paired seed bootstrap to LRBE-IG minus QualityFirst-IG count, LRBE-IG minus MinRisk predicted risk and Brier score, and the prespecified exposure-by-region-weight interactions. Its confirmation run holds training randomness fixed across branches. Candidate-label perturbations leave predictions, candidates, IG scores, and selections unchanged; test-label perturbations leave the fitted predictor unchanged.

For feedback--prediction associations, the primary variables are raw gains from the same-seed NoExplore reference: $\Delta F=F_m-F_N$, $\Delta\mathrm{Recall}=\mathrm{Recall}_m-\mathrm{Recall}_N$, and $\Delta\mathrm{Brier}=\mathrm{Brier}_N-\mathrm{Brier}_m$. Positive Brier gain means improvement. Spearman associations are computed separately within dataset and budget, retaining all methods belonging to a resampled seed in each of 10,000 cluster-bootstrap draws. Budget-centered associations are secondary descriptions, not causal estimates. Reference-gap ratios are supplemental only: $(M_m-M_N)/(M_O-M_N)$ is interpretable as recovery only when the oriented Oracle gap is positive; nonpositive or numerically unresolved gaps are N/A. The primary raw-gain analysis avoids dividing by these unstable reference differences.

\paragraph{Cohort provenance.}
UCI is the Default of Credit Card Clients dataset~\cite{yeh2016uci}; HMEQ is the home-equity cohort documented in~\cite{baesens2016creditrisk}. LendingClub merges the documented DePaul-hosted 2014 and 2015 public samples into a mature Fully Paid/Charged Off cohort~\cite{depaul2015lendingclub}. Its outcomes concern historically issued loans, not observed outcomes for genuinely rejected applicants. PPDai is linked to Figshare file 26049578~\cite{figshare26049578}; the current reconstructed source is recorded by hash, but byte parity with the historical processed file is not established. It therefore carries exploratory evidentiary weight and is not pooled with historical PPDai results. Its exposure column is identically one. Hash-checked source instructions identify the exact inputs without redistributing third-party data. The strict exposure audit finds no missing, nonnumeric, nonfinite or nonpositive exposure in the three stored exposure columns; UCI constructs unit exposure directly.

\section{Results}
\label{sec:results}

\subsection{How much quantity does score-first selection sacrifice?}
The first question is how many feasible labels a score-first policy gives up. Removing the count equality quantifies that opportunity cost. Relative to ScoreFirst-FAVE, LRBE-FAVE acquires 67.1 [62.6, 71.4] additional labels on HMEQ, 209.6 [202.5, 216.1] on LendingClub, 2.0 [0.9, 3.1] on PPDai, and 56.4 [54.7, 58.2] on UCI. Recall favors LRBE-FAVE on HMEQ by 0.0175 [0.0041, 0.0297]; the other three Recall intervals contain zero. The count constraint guarantees maximum cardinality within each policy's current state. The cumulative differences reported here are empirical outcomes of the resulting policy trajectories, and score-first selection does not produce a consistent predictive improvement across datasets.

The random-order control provides a second view. Against MinRisk, it acquires 87.0 fewer labels on HMEQ, 317.6 fewer on LendingClub, and 59.5 fewer on UCI; PPDai differs by 0.3. HMEQ Recall is lower by 0.0268 [-0.0395, -0.0140], while the other Recall intervals contain zero. The fixed-count random-score control remains within 2.7 labels of MinRisk in every dataset, and all Recall and Brier intervals contain zero. These are complete-policy comparisons over policy-specific trajectories.

FAVE and MinRisk have close observed cumulative counts, an empirical similarity across their policy-specific trajectories. FAVE's count differences are 1.1 [-0.9, 3.0] on HMEQ, $-0.1$ [-2.5, 2.6] on LendingClub, 0.0 [-1.0, 1.1] on PPDai, and $-0.3$ [-2.0, 1.5] on UCI. It uses more predicted risk in all four datasets, while every Recall interval contains zero. Three Brier intervals contain zero; UCI favors FAVE by $-0.00018$ [$-0.00028$, $-0.00008$]. The UCI estimate is a dataset-specific benefit amid otherwise uncertain Brier differences.

\subsection{Independent information gain separates quantity from score choice}
The IG extension tests the quantity constraint with a score independent of FAVE's linearization. All five policies are rerun together at $\rho=0.1$. Panel A of Table~\ref{tab:ig-primary} shows sizeable feedback increases for LRBE-IG on HMEQ, LendingClub, and UCI. Recall estimates span both directions, while UCI Brier error increases by 0.000156 [0.000019, 0.000298]. PPDai remains a small-feedback exploratory case.

\begin{table}[H]
\caption{Primary IG contrasts at $\rho=0.1$ ($n=10$ paired seeds). Entries are mean left-minus-right differences with pointwise 95\% intervals. Positive $\Delta$Brier denotes higher error.}
\label{tab:ig-primary}
\centering\scriptsize
\resizebox{\linewidth}{!}{%
\begin{tabular}{lrrrr}
\toprule
Dataset & $\Delta K$ & $\Delta\widehat R$ & $\Delta$Recall & $\Delta$Brier\\
\midrule
\multicolumn{5}{l}{\emph{Panel A: LRBE-IG $-$ ScoreFirst-IG}}\\
HMEQ & 31.3 [28.3,34.5] & 5,267.5 [$-$6,697.3,16,208.6] & .00535 [$-$.00068,.01170] & $-$.000302 [$-$.000937,.000364]\\
LendingClub & 49.5 [46.0,53.1] & $-$9,907.0 [$-$21,199.7,$-$1,269.3] & $-$.00073 [$-$.00524,.00425] & $-$.000070 [$-$.000251,.000103]\\
UCI & 40.1 [36.8,43.6] & 1.624 [.241,2.869] & .00166 [$-$.00187,.00517] & .000156 [.000019,.000298]\\
PPDai$^\dagger$ & 0.2 [0.0,0.5] & .050 [.000,.132] & $-$.00056 [$-$.00139,.00000] & $-$.000029 [$-$.000066,.000000]\\
\midrule
\multicolumn{5}{l}{\emph{Panel B: LRBE-IG $-$ CountOnly-MinRisk}}\\
HMEQ & 0.2 [$-$1.6,1.7] & 28,765.5 [19,361.2,39,425.1] & $-$.00518 [$-$.01423,.00388] & $-$.000068 [$-$.000479,.000292]\\
LendingClub & $-$2.9 [$-$6.0,0.1] & 27,025.7 [19,238.3,35,010.3] & .00045 [$-$.00286,.00488] & $-$.000029 [$-$.000174,.000120]\\
UCI & $-$1.2 [$-$3.0,0.5] & 1.827 [1.004,2.733] & $-$.00025 [$-$.00270,.00222] & $-$.000098 [$-$.000262,.000059]\\
PPDai$^\dagger$ & 0.4 [$-$0.6,1.5] & 1.024 [.497,1.606] & $-$.00052 [$-$.00270,.00142] & $-$.000064 [$-$.000134,.000005]\\
\bottomrule
\end{tabular}}
\vspace{1mm}
\parbox{\linewidth}{\scriptsize $^\dagger$The registered reconstructed PPDai artifact provides unit-exposure exploratory evidence and is analyzed separately from the historical PPDai block. Risk units are dataset-specific.}
\end{table}

Panel B shows that LRBE-IG and MinRisk have close cumulative counts, while LRBE-IG spends more predicted risk on every dataset. The Recall and Brier estimates span improvement and deterioration. In this block, IG changes composition and risk use more clearly than it changes downstream prediction.

The same-state audit fixes the MinRisk visible set, candidates, costs, and numerical budget. Across 100 states per dataset, LRBE-IG has zero count sacrifice. ScoreFirst-IG gives up 3.46, 5.37, 4.04, and 0.02 applicants per state on HMEQ, LendingClub, UCI, and PPDai, respectively. Its mean Jaccard overlap with MinRisk is 0.436, 0.602, 0.135, and 0.224, confirming that the score changes both cardinality and composition at the anchor state. Alternative sets do not generate later audit states, so the audit ends at the selection outcome.

\subsection{Controlled conditions reproduce the selection trade-off}
The one-step experiment holds the acquisition state and training randomness fixed while varying exposure heterogeneity and the evaluation weight of an underrepresented region. LRBE-IG selects 2.5 [1.45, 3.65] more applicants than ScoreFirst-IG under unit exposure and 3.2 [1.75, 4.85] more under heterogeneous exposure. Relative to MinRisk, it uses 0.221 [0.165, 0.284] and 0.136 [0.101, 0.171] more predicted-risk units, respectively. All four primary Brier intervals span zero.

The prespecified condition interactions are also uncertain. Increasing the target-region weight from 0.1 to 0.5 changes the IG-minus-MinRisk Brier contrast by 0.000126 [$-0.000005$, 0.000255] under unit exposure and $-0.000058$ [$-0.000217$, 0.000102] under heterogeneous exposure; the difference between those changes is $-0.000184$ [$-0.000413$, 0.000029]. These estimates leave the moderating roles of exposure and target weight unresolved. The incoming batch averages about 3,751 ordinary approvals, compared with exploratory $K^*$ means of 29.55 and 118.7, so the evaluated retraining intervention is small relative to the shared ordinary feedback.

\subsection{Budget sensitivity}
The budget ratio has a much larger and more consistent relation with acquired quantity than the secondary score. HMEQ MinRisk grows from 57.4 labels at $\rho=0.025$ to 374.0 at $\rho=0.4$; FAVE grows from 58.0 to 375.8. LendingClub MinRisk grows from 207.6 to 1,253.1 and FAVE from 206.3 to 1,253.3. Recall moves less and method ordering changes across budgets. Figure~\ref{fig:budgetsweep} shows representative HMEQ and LendingClub trajectories. The sweep changes both the scale of feedback and the subsequent policy state, so it is a sensitivity analysis rather than a quantity-only intervention.

\begin{figure}[t]
\centering
\paperfigure{fig_budget_sweep_feedback_recall.png}{\linewidth}
\caption{Mean cumulative exploratory count and final Recall across budget ratios for HMEQ and LendingClub (ten fresh seeds). Larger budgets increase feedback; the Recall ordering of MinRisk and FAVE varies across budgets.}
\label{fig:budgetsweep}
\end{figure}

\subsection{Feedback diagnostics vary across datasets and outcomes}
Normalized recovery divides policy changes by an Oracle--NoExplore gap and represents recovery when that oriented gap is positive. In the present data, the Recall gap is nonpositive in 10 of 40 dataset--seed references: two LendingClub seeds, two PPDai seeds, including one zero, and six UCI seeds. Oracle gaps are also nonpositive in 21 of 40 ROC-AUC and 15 of 40 PR-AUC references. We therefore use raw gains from NoExplore as the primary analysis and reserve normalized recovery for positive-gap cells.

Figure~\ref{fig:rawgain} relates visible-set default-ratio gain to Recall and Brier improvement at $\rho=0.1$. HMEQ has a positive Recall association, $\rho_s=0.397$ [0.013, 0.660], but its Brier association is uncertain, 0.062 [$-0.175$, 0.381]. LendingClub, PPDai, and UCI Recall intervals contain zero at this budget, as do all four Brier intervals. Across the five within-budget analyses, HMEQ Recall estimates range from 0.389 to 0.514 and all intervals exclude zero. LendingClub instead has negative Recall associations at $\rho=0.025$ and 0.05, followed by uncertain estimates. PPDai's smallest-budget cells have near-degenerate ranks and yield perfect boundary coefficients.

\begin{figure}[t]
\centering
\paperfigure{fig_raw_gain_associations.pdf}{0.91\linewidth}
\caption{Within-dataset associations between visible-set default-ratio gain and downstream improvement at $\rho=0.1$. Points use 40 terminal policy cells (four policies, ten seeds); bars are pointwise 95\% intervals from 10,000 whole-seed bootstrap resamples. Positive Brier gain means lower error.}
\label{fig:rawgain}
\end{figure}

These associations concern a diagnostic of the complete visible set, which includes the initial set and ordinary approvals. A larger ratio can move toward or away from the population target, and policy-specific predictors make selected mean PD an unsuitable cross-policy identity measure. Feedback diagnostics align with Recall in some regimes and diverge in others. We report the correlations as descriptive relationships between complete-policy outcomes.

\subsection{Predicted feasibility and realized default exposure}
All non-Oracle selections obey the fitted point-cost constraint within audit tolerance. Across 104 equally sized method--dataset--budget groups, the mean rate at which realized defaulted exposure exceeds the policy's numerical predicted budget is 0.4798, with group rates ranging from 0.01 to 0.84. Because each group contains 100 seed--round observations, the unweighted group mean also equals the observation-weighted mean for these included cells. The range summarizes observed groups rather than sampling uncertainty.

This comparison evaluates an ex-post loss exceedance against a budget constructed from fitted expected costs rather than a realized-loss quantile. The exceedance rate shows the gap between point-estimate feasibility and ex-post loss-cap compliance, combining model calibration, stochastic default variation, exposure assumptions, and dependence among loans.

\paragraph{Descriptive computation time.}
Recorded episode wall time includes training, selection, evaluation, and system effects. Mean LendingClub times in the IG block are 44.0 s for MinRisk, 82.8 s for LRBE-FAVE, and 498.6 s for LRBE-IG; these are end-to-end episode descriptions rather than isolated acquisition latency.

\section{Discussion and Scope}
\label{sec:discussion}

\paragraph{What a quantity-preserving reference contributes.}
Maximizing feasible count makes the opportunity cost of a secondary acquisition score observable. The scientific question begins after that arithmetic fact: does the score earn the observations it displaces? FAVE, IG, and the random-order control all produce sizeable quantity differences when the maximum-count equality is removed, yet composition-first selection has no consistent predictive advantage. Conversely, the additional labels preserved by LRBE do not reliably improve prediction; UCI Brier is slightly worse for LRBE-IG than for QualityFirst-IG. The useful outcome is therefore a reference comparison, not a universal ordering. Within the eligibility rule and predicted-risk constraint, MinRisk states the least predicted risk required to retain the current maximum count.

\paragraph{Policy effects and composition evidence.}
The three experimental views answer different questions. Dynamic replay compares complete acquisition policies that begin with the same data stream but evolve through policy-specific feedback, predictions, candidates, and budgets. Same-state audits fix those quantities and show exactly how selectors exchange count, risk, and score at one decision point. The controlled experiment adds a single retraining step after fixing the acquisition state. It reproduces the count and predicted-risk contrasts, while its exposure and target-weight interactions remain uncertain. Taken together, the results locate the trade-off at selection but do not reduce long-run policy differences to composition alone.

\paragraph{Why predictive differences can remain small.}
Most labels added in the controlled task come from ordinary approvals shared across selectors: about 3,751 per incoming batch, compared with exploratory $K^*$ means of 29.55 under unit exposure and 118.7 under heterogeneous exposure. Differences among exploratory selections therefore modify only a fraction of the post-acquisition training set. This scale contrast offers one plausible account of the small Brier differences, but the experiment was designed to test exposure and target-weight interactions rather than isolate dilution. More broadly, an acquisition score is useful only when the labels it changes affect the fitted predictor and the target metric; improved score totals or altered feedback composition are intermediate outcomes.

\paragraph{Scope of inference.}
The evidence concerns random-order replay on historical cohorts with observed outcomes. LendingClub does not contain outcomes for actual historical rejects, labels arrive immediately in the simulation, and $\tau=0.5$ is a common experimental operating point rather than a production cutoff. PPDai contributes unit-exposure exploratory evidence; monetary heterogeneity comes from HMEQ and LendingClub. The IG comparison uses the registered primary budget $\rho=0.1$, while the broader budget sweep uses the original policy set. Published online exploration methods motivate the question, but the randomized controls are not implementations of those algorithms. The conclusions apply to the evaluated acquisition policies and replay conditions.

\paragraph{Risk interpretation.}
An optimizer can enforce a sum of predicted losses without controlling the tail of realized loss. Here, an exceedance means that realized defaulted exposure in a seed--round observation is larger than the policy's point-estimate budget. It is an ex-post comparison with an expected-loss quantity, not a direct estimate of safety failure. The observed frequency does not attribute the difference separately to probability error, default randomness, or exposure assumptions. Live credit experimentation would require product-specific exposure and recovery definitions, delayed-outcome handling, concentration and fairness safeguards, and explicit stopping rules. None of these operational guarantees follows from the current point-cost constraint.

\section{Conclusion}
\label{sec:conclusion}
Risk-budgeted feedback acquisition involves two distinct choices: how many outcomes to make observable and which applicants should provide them. LRBE separates these choices by preserving the current maximum feasible count before applying a composition score. Across FAVE, independent IG, and randomized controls, removing that constraint often sacrifices substantial feedback, but preserving more labels does not itself guarantee better prediction. At nearly the same cumulative count, FAVE and IG also spend more predicted risk than MinRisk without a consistent downstream return. The controlled analysis reproduces these selection trade-offs but does not identify exposure heterogeneity or target-region emphasis as reliable moderators. Acquisition studies should therefore report quantity, predicted-risk use, composition change, and downstream performance together. Minimum-risk maximum-count selection provides the reference needed to decide whether a more elaborate score earns its cost.

\appendix
\section{Secondary Score and Reproducibility}
\label{app:fave}\label{app:repro}
\paragraph{Frozen FAVE-v2.1 score.}
Let $v_j=[1,x_j^\top]^\top$ be the augmented, preprocessed feature vector. For the current visible set $\mathcal L_t$ and label-free reference batch $\mathcal U_t$, define
\begin{align}
H_t&=\sum_{j\in\mathcal L_t}\max\{\widehat p_j(1-\widehat p_j),10^{-6}\}v_jv_j^\top,\\
G_t&=|\mathcal U_t|^{-1}\sum_{j\in\mathcal U_t}\widehat p_jv_jv_j^\top,\qquad
\eta=\max\{10^{-6}\operatorname{tr}(H_t)/d,10^{-10}\},\\
A_t&=H_t+\eta I,\qquad
\widetilde w_i=\max\{\widehat p_i(1-\widehat p_i)\widehat p_i^{0.5},10^{-10}\},\\
q_i^{\mathrm{FAVE}}&=
\frac{\widetilde w_i v_i^\top A_t^{-1}G_t A_t^{-1}v_i}
{1+\widetilde w_i v_i^\top A_t^{-1}v_i}.
\end{align}
Here $d$ includes the intercept coordinate. Predicted probabilities are clipped to $[10^{-9},1-10^{-9}]$ before these operations; nonfinite score outputs are zeroed and scores are clipped below at zero. The default-weighting exponents and numerical regularization are fixed. This linearized feature-space construction is used with boosted-tree predictions as a scoring proxy; it is not an exact influence calculation for the tree-fitting procedure. The Joint comparator selects $\lambda\in\{0.01,0.1,0.5,1,2,5,10\}$ using only initial-set validation. The initial split is stratified where possible, with a deterministic random split as fallback.

\paragraph{Input and trajectory audit.}
Exact-input SHA-256 checks cover HMEQ, LendingClub and PPDai. All 78,122 records have numeric, finite, positive exposure; PPDai values are all one. UCI constructs ones directly. The strict input loader rejects invalid exposure rather than replacing it by zero. Eight validation trajectories (four datasets, seed 1001, MinRisk and FAVE) agree with the archived outputs in the checked counts, risk values, and predictive metrics. This check is not an additional statistical replication. Of 40 paired post-initial rounds, incoming-batch hashes agree in all 40, while candidates and numerical budgets agree in only four. This directly documents the policy-specific state evolution.
The IG and controlled-condition archives retain their own registered seeds and are not pooled as independent replications. The controlled confirmation uses common branch randomness and label-perturbation checks; all 320 endpoint rows match the original block in selection, risk, Recall, and Brier score. Source hashes, full-precision tables, parameter grids, and executable tests accompany the archive.

\bibliographystyle{splncs04}
\bibliography{../literature/refs,../literature/dataset_refs,../literature/closest_work}
\end{document}
